% GENERATED FILE. Edit canonical lessons, then run npm run generate:books.
% canonical-source-hash: fnv1a64:11b4ce013501604e
% canonical-lessons: GU-C196-makan, GU-C196-mal-storey, GU-C196-lipht, GU-C196-pesvu, GU-C196-praves

\chapter{Up to the top storey}
\label{ch:gu-up-to-the-top-storey}
\hlchaptermodality{196}

\begin{chapteropening}
\textbf{By the end of this chapter:} \emph{I can find my way in a building: the floor, the lift and the way in.}

Complete the last lesson of chapter 196: I can find my way in a building: the floor, the lift and the way in.
\end{chapteropening}

\section[\texorpdfstring{\gu{મકાન} (makān)}{મકાન (makān)}]{\gu{મકાન} (makān) — building, house}

\label{lesson:GU-C196-makan}



\begin{quote}
Before the new one: say the Gujarati for whereabouts, then the Gujarati for distance.
\end{quote}

\subsection*{You'll want to know: \gu{મકાન}}
\textbf{\gu{મકાન}} — \emph{makān} — \textquotedblleft{}building, house\textquotedblright{}.

\begin{quote}
\textbf{\gu{આ} \gu{મકાન} \gu{જૂનું} \gu{છે}.} — \emph{ā makān jūnũ che} — \textquotedblleft{}This building is old.\textquotedblright{}
\end{quote}

\begin{cousinweb}
\textbf{\gu{મકાન}} (\emph{makān}) is a loan from Arabic \emph{makān}, \textquotedblleft{}place\textquotedblright{}, which reached Gujarati through Persian. Its Arabic root is \emph{k-w-n}, \textquotedblleft{}to be\textquotedblright{} (Lane).
\end{cousinweb}

\begin{grammarlens}[title={What you've built}]
One more word for everyday life at A2.
\end{grammarlens}

\subsection*{Your turn}
\begin{itemize}
\raggedright
  \item \emph{Say it:} \emph{makān}
  \item \emph{Say it:} \emph{makān}, once more
  \item \emph{Say it:} \emph{ā makān jūnũ che} — This building is old.
  \item \emph{Recall:} say the Gujarati for whereabouts, then the Gujarati for distance, then say \emph{makān} again
\end{itemize}

\begin{tcolorbox}[breakable,colback=teal!4,colframe=teal!35!black,title={Before you move on}]
What does \gu{મકાન} mean? (\textquotedblleft{}Building, house\textquotedblright{}.) Say it once more.
\end{tcolorbox}
\section[\texorpdfstring{\gu{માળ} (māḷ)}{માળ (māḷ)}]{\gu{માળ} (māḷ) — storey, floor}

\label{lesson:GU-C196-mal-storey}



\begin{quote}
Before the new one: say the Gujarati for distance, then the Gujarati for building.
\end{quote}

\subsection*{You'll want to know: \gu{માળ}}
\textbf{\gu{માળ}} — \emph{māḷ} — \textquotedblleft{}storey, floor\textquotedblright{}.

\begin{quote}
\textbf{\gu{અમારું} \gu{ઘર} \gu{ત્રીજા} \gu{માળે} \gu{છે}.} — \emph{amārũ ghar trījā māḷe che} — \textquotedblleft{}Our home is on the third floor.\textquotedblright{}
\end{quote}

A masculine noun.

\begin{cousinweb}
\textbf{\gu{માળ}} (\emph{māḷ}) is listed by CDIAL, glossed \textquotedblleft{}loft, upper storey\textquotedblright{}, under \emph{*māḍa}, \textquotedblleft{}upper storey of a house\textquotedblright{}, a reconstructed form that it takes from Dravidian, as in Tamil \emph{māṭam}, \textquotedblleft{}storeyed house\textquotedblright{} (CDIAL 9996).
\end{cousinweb}

\begin{grammarlens}[title={What you've built}]
One more word for everyday life at A2.
\end{grammarlens}

\subsection*{Your turn}
\begin{itemize}
\raggedright
  \item \emph{Say it:} \emph{māḷ}
  \item \emph{Say it:} \emph{māḷ}, once more
  \item \emph{Say it:} \emph{amārũ ghar trījā māḷe che} — Our home is on the third floor.
  \item \emph{Recall:} say the Gujarati for distance, then the Gujarati for building, then say \emph{māḷ} again
\end{itemize}

\begin{tcolorbox}[breakable,colback=teal!4,colframe=teal!35!black,title={Before you move on}]
What does \gu{માળ} mean? (\textquotedblleft{}Storey, floor\textquotedblright{}.) Say it once more.
\end{tcolorbox}
\section[\texorpdfstring{\gu{લિફ્ટ} (lipht)}{લિફ્ટ (lipht)}]{\gu{લિફ્ટ} (lipht) — elevator, lift}

\label{lesson:GU-C196-lipht}



\begin{quote}
Before the new one: say the Gujarati for building, then the Gujarati for storey.
\end{quote}

\subsection*{You'll want to know: \gu{લિફ્ટ}}
\textbf{\gu{લિફ્ટ}} — \emph{lipht} — \textquotedblleft{}elevator, lift\textquotedblright{}.

\begin{quote}
\textbf{\gu{આ} \gu{મકાનમાં} \gu{લિફ્ટ} \gu{નથી}.} — \emph{ā makānmā̃ lipht nathī} — \textquotedblleft{}This building has no lift.\textquotedblright{}
\end{quote}

A feminine noun.

\begin{cousinweb}
\textbf{\gu{લિફ્ટ}} (\emph{lipht}) is a loan from English \emph{lift}.
\end{cousinweb}

\begin{grammarlens}[title={What you've built}]
One more word for everyday life at A2.
\end{grammarlens}

\subsection*{Your turn}
\begin{itemize}
\raggedright
  \item \emph{Say it:} \emph{lipht}
  \item \emph{Say it:} \emph{lipht}, once more
  \item \emph{Say it:} \emph{ā makānmā̃ lipht nathī} — This building has no lift.
  \item \emph{Recall:} say the Gujarati for building, then the Gujarati for storey, then say \emph{lipht} again
\end{itemize}

\begin{tcolorbox}[breakable,colback=teal!4,colframe=teal!35!black,title={Before you move on}]
What does \gu{લિફ્ટ} mean? (\textquotedblleft{}Elevator, lift\textquotedblright{}.) Say it once more.
\end{tcolorbox}
\section[\texorpdfstring{\gu{પેસવું} (pesvũ)}{પેસવું (pesvũ)}]{\gu{પેસવું} (pesvũ) — to go in, to enter}

\label{lesson:GU-C196-pesvu}



\begin{quote}
Before the new one: say the Gujarati for storey, then the Gujarati for elevator.
\end{quote}

\subsection*{You'll want to know: \gu{પેસવું}}
\textbf{\gu{પેસવું}} — \emph{pesvũ} — \textquotedblleft{}to go in, to enter\textquotedblright{}.

\begin{quote}
\textbf{\gu{તે} \gu{મકાનમાં} \gu{પેઠો}.} — \emph{te makānmā̃ peṭho} — \textquotedblleft{}He went into the building.\textquotedblright{}
\end{quote}

Past \textbf{\gu{તે} \gu{પેઠો}}, \emph{te peṭho}.

\begin{cousinweb}
\textbf{\gu{પેસવું}} (\emph{pesvũ}) is a tadbhava, a word inherited from Sanskrit and reshaped by sound change. It comes from \emph{práviśati}, \textquotedblleft{}enters\textquotedblright{}, and CDIAL glosses the Gujarati verb \textquotedblleft{}enters\textquotedblright{} (CDIAL 8803).
\end{cousinweb}

\begin{grammarlens}[title={What you've built}]
One more word for everyday life at A2.
\end{grammarlens}

\subsection*{Your turn}
\begin{itemize}
\raggedright
  \item \emph{Say it:} \emph{pesvũ}
  \item \emph{Say it:} \emph{pesvũ}, once more
  \item \emph{Say it:} \emph{te makānmā̃ peṭho} — He went into the building.
  \item \emph{Recall:} say the Gujarati for storey, then the Gujarati for elevator, then say \emph{pesvũ} again
\end{itemize}

\begin{tcolorbox}[breakable,colback=teal!4,colframe=teal!35!black,title={Before you move on}]
What does \gu{પેસવું} mean? (\textquotedblleft{}To go in, to enter\textquotedblright{}.) Say it once more.
\end{tcolorbox}
\section[\texorpdfstring{\gu{પ્રવેશ} (praveś)}{પ્રવેશ (praveś)}]{\gu{પ્રવેશ} (praveś) — entry, way in}

\label{lesson:GU-C196-praves}



\begin{quote}
Before the new one: say the Gujarati for elevator, then the Gujarati for to go in.
\end{quote}

\subsection*{You'll want to know: \gu{પ્રવેશ}}
\textbf{\gu{પ્રવેશ}} — \emph{praveś} — \textquotedblleft{}entry, way in\textquotedblright{}.

\begin{quote}
\textbf{\gu{પ્રવેશ} \gu{પાછળથી} \gu{છે}.} — \emph{praveś pāchaḷthī che} — \textquotedblleft{}The way in is at the back.\textquotedblright{}
\end{quote}

A masculine noun.

\begin{cousinweb}
\textbf{\gu{પ્રવેશ}} (\emph{praveś}) is Sanskrit \emph{praveśa}, \textquotedblleft{}entering, entrance\textquotedblright{}, taken over whole as a learned loan, a tatsama. It shares its root with \emph{práviśati}, \textquotedblleft{}enters\textquotedblright{} (CDIAL 8803).
\end{cousinweb}

\begin{grammarlens}[title={What you've built}]
One more word for everyday life at A2.
\end{grammarlens}

\subsection*{Your turn}
\begin{itemize}
\raggedright
  \item \emph{Say it:} \emph{praveś}
  \item \emph{Say it:} \emph{praveś}, once more
  \item \emph{Say it:} \emph{praveś pāchaḷthī che} — The way in is at the back.
  \item \emph{Recall:} say the Gujarati for elevator, then the Gujarati for to go in, then say \emph{praveś} again
\end{itemize}

\begin{tcolorbox}[breakable,colback=teal!4,colframe=teal!35!black,title={Before you move on}]
What does \gu{પ્રવેશ} mean? (\textquotedblleft{}Entry, way in\textquotedblright{}.) Say it once more.
\end{tcolorbox}
